\section{Where the Full Decoder Helps}
\label{sec:study}

Our goal is to characterize \emph{where} the extra computation in a full 3D
decoder changes predictions relative to an efficient decoder, rather than to
summarize it with a single accuracy or FLOPs number.  This section defines the
matched comparison, the per-voxel prediction transitions we count, the
test-time signal we evaluate, and the region-level opportunity analysis.

\paragraph{Matched decoder comparison.}
We compare two models that share an identical encoder and differ only in the
decoder.  \fullmodel{} uses the original wide, high-resolution decoder;
\effimodel{} uses the \method{} decoder, which uniformly reduces channel width
and removes the highest-resolution stages~\cite{rahman2025effidec3d}.  Because
the encoder is held fixed, any difference in the two models' predictions is
attributable to decoder capacity alone; the study therefore says nothing about
the encoder, which we leave untouched.  Let $\hat{y}_{\mathrm{E}}(v)$ and
$\hat{y}_{\mathrm{F}}(v)$ denote the efficient and full predicted labels at
voxel $v$, and $y(v)$ the ground-truth label.

\paragraph{Prediction transitions.}
Rather than treating every disagreement between the two decoders as an
improvement, we separate the two directions of change with plain indicator
counts.  A \emph{positive transition} is a voxel the efficient decoder gets
wrong and the full decoder gets right; a \emph{negative transition} is the
opposite:
\begin{align}
P(v) &= \mathbb{1}\!\left[\hat{y}_{\mathrm{E}}(v)\neq y(v)\ \wedge\ \hat{y}_{\mathrm{F}}(v)=y(v)\right],\\
N(v) &= \mathbb{1}\!\left[\hat{y}_{\mathrm{E}}(v)=y(v)\ \wedge\ \hat{y}_{\mathrm{F}}(v)\neq y(v)\right],\\
U(v) &= P(v)-N(v).
\label{eq:transitions}
\end{align}
$U(v)$ is the \emph{net} transition.  The corresponding rates over a volume are
$R_{\mathrm{pos}}=\mathrm{mean}_v P(v)$, $R_{\mathrm{neg}}=\mathrm{mean}_v N(v)$,
and $R_{\mathrm{net}}=R_{\mathrm{pos}}-R_{\mathrm{neg}}$.  Reporting $P$, $N$,
and $U$ together prevents the common error of counting negative transitions as
part of the decoder's benefit.

\paragraph{Test-time signal.}
A benefit that could guide computation must be identifiable \emph{without} the
ground truth.  We therefore evaluate whether the efficient model's own
uncertainty predicts where net transitions occur.  The primary signal is the
per-voxel predictive entropy of the efficient model,
\begin{equation}
H(v) = -\sum_{c=1}^{C} p_c(v)\,\log p_c(v),
\label{eq:entropy}
\end{equation}
where $p_c(v)$ is its softmax probability for class $c$ and $C$ is the number of
classes.  We compare $H$ against confidence $\max_c p_c(v)$ and Monte-Carlo
dropout~\cite{gal2016dropout} as alternative cheap signals.

\paragraph{Region-level opportunity analysis.}
Isolated voxels cannot be processed efficiently by standard 3D convolutions, so
a deployable analysis must operate on contiguous regions.  We partition the
foreground-union volume into non-overlapping $16^3$ blocks, rank blocks by mean
entropy, and select the top fraction of blocks at a predeclared budget.  Each
selected block is expanded by a $4$-voxel context halo, and we report the
\emph{executed} volume fraction this induces.  Within the selected region we
measure net transition utility, defined as $(\sum P - \sum N)/\sum P$ over the
foreground union, and compare it against a matched-random selector that draws
the same number of blocks.  A voxel-level ranking of the same signal is reported
only as a non-deployable oracle upper bound.

\paragraph{Statistics.}
All confidence intervals resample \emph{subjects}, not voxels: with only twelve
validation volumes, treating millions of correlated voxels as independent would
grossly overstate significance.  Correlations between $H$ and $U$ are computed
per subject over entropy quantile bins and aggregated across subjects; the
opportunity comparison uses a paired subject bootstrap in which the entropy and
random selectors share the same resampled subject indices.
